\section{Experiments}
\label{sec:experiments}

The experiments test one sealed score in two distinct roles. RQ1 asks whether
it prospectively orders later damage. RQ2 tests, on each setting's declared
fidelity support, whether its forward-only coordinate recovers the ordering
induced by exact Loss Susceptibility. RQ3 asks whether
the same score improves a real allocation decision under a fixed repair
budget. Table~\ref{tab:pred-value} decides RQ1,
Table~\ref{tab:loss-susceptibility-fidelity} decides RQ2, and
Table~\ref{tab:prot-effect} decides RQ3;
Table~\ref{tab:robustness} reports coverage and failure boundaries.

% Claim-bearing minimal layout for RQ1--RQ3. The four tables are deliberately
% spread across the subsections they decide (I before 5.1, II after 5.2,
% III after 5.3, IV after 5.4) rather than clustered in one include.
% Fill from the evidence pipeline, never by hand:
%   [ledger]   results/paper/evidence_ledger.json
%   [derived]  derivable from ledger fields
%   [cert]     runs/<setting>/fidelity/<model>.json (registered certificate)
%   [export+]  needs a small export extension
% NOTE: sections/generated/table_core_evidence.tex and its renderer
% (src/rsus/evidence/tables.py via experiments/paper/build_evidence.py)
% still emit the previous five-table layout and must be re-synced to this one.


\subsection{Sealed Campaign Design and Evidence Rules}
\label{sec:benchmark-protocol}

\paragraph{Sealed data roles.}
Candidates are split by semantic group into discovery and audit folds.
Three disjoint, target-request-disjoint development folds separately fix the
operating point $(R,\eta)$ for the Loss Susceptibility estimator
and fidelity fraction $q_{\mathrm{fid}}$, the shared mixture weight
$\widehat\alpha_{\mathrm{pred}}$, and the repair protocol. The latter includes
$K_p$, $B_{\mathrm{tok}}$, stopping rules, and guards. On a target
request, discovery candidates may fit rank maps and supply repair examples;
audit candidates never enter repair. No target damage or audit outcome can
change the score, pool, parent checkpoint, or returned repair checkpoint.

\paragraph{Settings and parents.}
The executed campaign spans TOFU \cite{maini2024tofu} on Qwen2.5-1.5B,
Qwen2.5-7B, Qwen2.5-14B, and Llama-3.1-8B, and WMDP-bio with an MMLU retain
universe \cite{li2024wmdp,hendrycks2021measuring} on Qwen2.5-7B and
Qwen2.5-14B. The claims carry deliberately different scopes. The
setting-level estimator claim is certified per model at each model's own
frozen operating point, so it spans the model axis. The prediction and
repair claims are carried by TOFU on Qwen2.5-7B, the primary setting, which
holds the complete sealed chain: fidelity certificate, frozen objective
settings, prospective audit, mixture development, and fixed-budget repair.
The remaining settings are predeclared replications whose later phases are
reported as coverage in Table~\ref{tab:robustness} rather than as
additional confirmatory evidence. Each
setting's core parent roster is reduced from the full seven-objective grid
(GradDiff, NPO, SimNPO, GRU, RMU, RepNoise, Circuit Breakers, with GA and
IdkDPO as stress controls) by development calibration alone, before any
audit or probe outcome exists; demoted objectives keep best-effort stress
settings so every column is still measured. Per-setting forgetting and
utility bounds are fixed by dated pre-audit amendments
(\S\ref{sec:limitations}). A campaign manifest freezes every
model--dataset--request--parent cell and seed.

\paragraph{Prediction references.}
The deployed score $S_{\widehat\alpha_{\mathrm{pred}}}$ is compared with its
two endpoints: $S_0$ uses the rank-normalized Loss Susceptibility
estimate only, and $S_1$ uses Representation Proximity only. Additional references include exact
Loss Susceptibility, checkpoint-local candidate--forget-gradient
alignment computed at $\theta_0$, initial NLL, answer length, lexical and
sentence-embedding proximity, request-shuffled
Representation Proximity, a one-direction loss response, and repeated random
rankings. The strongest simple control is chosen on development requests and
frozen before target evaluation.

\paragraph{Protection comparators.}
The two-signal selector, the two endpoint selectors, and repeated random
selection use exactly the same repair harness. The unrepaired entry checkpoint
measures net repair benefit. The $q_G^\star$-substituted mixture selector,
uniform rehearsal, and a
realized-discovery-damage oracle are diagnostic arms that separate estimator
error, repair capacity, and selector quality; they do not replace the
confirmatory comparisons.
Retain-aware mitigation methods that reweight or curate the forget data,
constrain the parent objective, or rectify its gradients intervene on the
parent trajectory itself; running them as arms would confound selector quality
with the operator and objective differences this harness deliberately holds
fixed. They are therefore contrasted in the related-work discussion rather
than included as confirmatory comparators.

\paragraph{Outcomes and aggregation.}
The primary retained-behavior outcome is candidate-level $\Delta$NLL relative
to $\theta_0$. We report request-level mean damage and
$\mathrm{CVaR}_{.95}$, the mean of the worst 5\% of candidate damage, together
with each dataset's native retained metric and forgetting/utility slack.
Prediction associations are Spearman rank correlations computed within each
request--seed--parent cell. Seeds are averaged within requests and requests
receive equal weight. A paired hierarchical bootstrap resamples requests,
trajectory seeds, and semantic groups.

\paragraph{How the three claims are decided.}
RQ1 requires a positive prospective rank association, paired improvement over
both signal endpoints and the strongest simple control, and above-chance
recovery of the harmful tail on sufficient support; the rank and tail
conditions are decided jointly in Table~\ref{tab:pred-value}. For compact
table notation, $g_G$ is the gain over $S_0$, $g_H$ is the gain over $S_1$,
and $g_{\mathrm{ctl}}$ is the gain over the frozen simple control. The table
additionally reports the frozen mixture weight
$\widehat\alpha_{\mathrm{pred}}$ of every cell; an unresolved fallback weight
is claim-ineligible and marked~$\dagger$.

RQ2 is a setting-level numerical-estimator claim. It uses rank agreement
$f_\rho$ between estimated and exact Loss Susceptibility,
$\widehat q_G$ and $q_G^\star$, and their Top-$K_{\mathrm{fid}}$ overlap $f_K$
on the setting's declared fidelity support. The one-sided bounds for $f_\rho$
and $f_K$ resample that support. Each target profile separately uses
split-bank rank agreement, perturbation survival, and integrity coverage as
outcome-blind eligibility checks; these checks are not an exact-rank
certificate for its candidate support. Time and memory are also reported.
Predictive usefulness belongs to RQ1 rather than RQ2.

RQ3 compares the two-signal repair arm with each single-signal arm, repeated
random selection, and no repair. It requires lower held-out mean and tail
damage, non-inferiority on the native retained metric, and complete common support
inside the same forgetting/utility feasibility region; the effect and
contract conditions are decided jointly in Table~\ref{tab:prot-effect}.
Eligibility and
statistical success are reported separately. Appendix~\ref{app:outcomes}
specifies thresholds, missingness, multiplicity, and the complete decision
rules.

\subsection{The Sealed Score Prospectively Orders Retained Damage}
\label{sec:exp-prediction}

\paragraph{Answer (RQ1).}
{\setlength{\emergencystretch}{3em}
\PredictionHeadline\ \TailHeadline\par}

\paragraph{Candidate-level ordering.}
We compare the score sealed at $\theta_0$ with damage revealed at
$\theta_{t^\dagger}$ on the untouched audit fold. On the primary setting,
the sealed mixture attains a positive within-request rank correlation with
realized damage for both core parents, with positive one-sided lower bounds
on all three paired gains ($g_G$, $g_H$, $g_{\mathrm{ctl}}$) on identical
candidate support (Table~\ref{tab:pred-value}, left block). Reporting the
mixture weight, the absolute association, and all three gains prevents a
strong endpoint or a trivial checkpoint feature from being hidden by the
combined score.

\paragraph{Harmful-tail recovery.}
We also test whether the most damaged behaviors are found early enough to
prioritize. We predeclare an audit-tail fraction
$q_{\mathrm{tail}}$ and compare the top
$K_{\mathrm{tail}}=\lceil q_{\mathrm{tail}}
|\mathcal C_{\mathrm{audit}}|\rceil$ sealed ranks with the
$K_{\mathrm{tail}}$ largest positive-damage outcomes. The sealed mixture
recovers this tail above its exact finite-sample chance level
$q^{\mathrm{eff}}_{\mathrm{tail}}$ on both core parents
(Table~\ref{tab:pred-value}, right block), with positive-damage support and
tail-eligibility coverage reported alongside. This prediction-tail
measure is distinct from the $\mathrm{CVaR}_{.95}$ outcome used to evaluate
repair in \S\ref{sec:exp-protection}.

\paragraph{Evidence and boundaries.}
Table~\ref{tab:pred-value} decides RQ1 on the primary setting; the
replicated settings enter through Table~\ref{tab:robustness}. Appendix
Table~\ref{tab:prospective-prediction} expands the demoted and stress-control
parents and the frozen pre-unlearning reference roster, together with the
extra model access each reference predictor costs. Reversals against an endpoint,
diffuse or unsupported tails, invalid profiles, and gate-non-reaching runs
remain visible rather than being averaged away.

% =============================================================================
% Table II -- RQ2: setting-level estimator fidelity at the frozen operating point.
% =============================================================================
\begin{table}[t]
  \caption{\textbf{Loss Susceptibility estimation fidelity across four models
  and two benchmarks (RQ2).} Brackets are margins over the registered
  thresholds $\tau_\rho=0.80$ and $\tau_K=0.70$ at the frozen operating point
  $(R=64,\eta=3\!\times\!10^{-3})$.}
  \label{tab:loss-susceptibility-fidelity}
  \centering
  \small
  \setlength{\tabcolsep}{5pt}
  \begin{tabular}{@{}llccl@{}}
    \toprule
    Benchmark & Model & $f_\rho$ [$-\tau_\rho$] & $f_K$ [$-\tau_K$] & Cert. \\
    \midrule
TOFU     & Qwen2.5-1.5B  & 0.903 [+0.077] & 0.923 [+0.223] & pass \\
         & Qwen2.5-7B    & 0.916 [+0.099] & 0.872 [+0.146] & pass \\
         & Qwen2.5-14B   & 0.917 [+0.117] & 0.769 [+0.069] & pass \\
         & Llama-3.1-8B  & 0.888 [+0.083] & 0.821 [+0.069] & pass \\
\addlinespace[2pt]
WMDP-bio & Qwen2.5-7B    & 0.975 [+0.175] & 0.700 [+0.000] & pass \\
         & Qwen2.5-14B   & 0.970 [+0.170] & 1.000 [+0.300] & pass \\
    \bottomrule
  \end{tabular}
\end{table}

\subsection{Forward-Only Estimation Recovers Exact Loss Susceptibility}
\label{sec:exp-fidelity}

\paragraph{Answer (RQ2).}
{\setlength{\emergencystretch}{3em}\FidelityHeadline\par}

\paragraph{Ranking fidelity.}
At the operating point frozen on the calibration fold we compare the
forward-only estimate $\widehat q_G$ with exact Loss Susceptibility
$q_G^\star$ on the declared setting-level support. Every setting passes:
rank agreement $f_\rho$ reaches $0.888$--$0.975$ and clears
$\tau_\rho=0.80$ by at least $+0.077$
(Table~\ref{tab:loss-susceptibility-fidelity}). Where a three-seed
direction-bank sweep exists, the reported margin is taken at the seed
minimum, so the gate holds for every draw rather than on average. The
agreement spans Qwen2.5 at 1.5B, 7B, and 14B and Llama-3.1-8B across two
benchmarks with disjoint retain universes, so it is not a property of one
architecture or one corpus.

\paragraph{The binding quantity.}
Top-$K_{\mathrm{fid}}$ overlap is the tighter test, and we read it rather
than $f_\rho$ as the limit of what a certificate licenses. It is computed on
$13$ of $128$ TOFU and $10$ of $104$ WMDP candidates, so it moves in steps of
$1/K_{\mathrm{fid}}$: WMDP-bio on Qwen2.5-7B recovers seven of its ten
screening candidates and therefore meets $\tau_K=0.70$ exactly, while the
remaining settings clear it by $+0.069$ to $+0.300$. A single displaced
candidate would fail that certificate, and we treat its screening claim
accordingly.

\paragraph{Access cost.}
The direct-gradient reference streams one candidate gradient at a time; the
estimator uses $2R$ batched pool-forward sweeps and no candidate-side
backward pass during target profiling. Probe efficiency is eff/$\eta=1$ and
the random-projection reference agrees with the exact target at
$\rho_{AB}\ge0.911$ in every setting, so the residual gap is estimator
error, not a probe artifact.

\paragraph{What this claim excludes.}
RQ2 is rank-based, at the frozen operating point, on the declared fidelity
support. It does not establish magnitude accuracy, does not exact-rank
certify a target candidate support, and does not show that Loss
Susceptibility predicts damage---that is the endpoint and control comparison
in \S\ref{sec:exp-prediction}. Appendix Table~\ref{tab:bwfree} gives the
nested-$R$ fidelity--cost frontier and integrity checks; it separates the two
quantities, with rank agreement still improving at $R=128$ while
top-$K_{\mathrm{fid}}$ overlap does not improve beyond the frozen operating
point.

% =============================================================================
% Table III -- RQ3 (effect + contract, decided jointly): fixed-budget repair
%              on the primary setting.
% =============================================================================
\begin{table*}[t]
  \caption{\textbf{Fixed-budget protection (RQ3): allocation is the only free
  variable.} Setting: TOFU, Qwen2.5-7B. All arms share the repair operator,
  selection quota $K_p$, token budget, guards, and feasibility contract;
  only the selector differs. Damage cells report mean\,;\,$\mathrm{CVaR}_{.95}$
  of held-out audit $\Delta$NLL (lower is better). The contract block reports
  the minimum forgetting and native-utility slack across arms and the number
  of feasible confirmatory arms; RQ3 (\texttt{E/P}) passes only when the
  two-signal arm improves mean and tail damage over every same-budget
  comparator and the no-repair checkpoint with the contract intact.}
  \label{tab:prot-effect}
  \centering
  \small
  \setlength{\tabcolsep}{4pt}
  \begin{tabular}{@{}lccccccccc@{}}
    \toprule
    & \multicolumn{4}{c}{Same-budget comparators}
      & & \multicolumn{3}{c}{Contract} & \\
    \cmidrule(lr){2-5}\cmidrule(lr){7-9}
    Parent
      & No repair                                  % [ledger] entry checkpoint
      & Random                                     % [ledger] repeated draws
      & $S_0$                                      % [ledger]
      & $S_1$ (prox.)                              % [ledger]
      & $S_{\widehat\alpha_{\mathrm{pred}}}$ (ours)% [ledger]
      & Forget slack                               % [ledger] min over arms
      & Utility slack                              % [ledger] native non-inferiority
      & Feasible arms                              % [ledger]
      & RQ3 E/P \\
    \multicolumn{10}{@{}r}{\footnotesize all damage cells: mean\,;\,$\mathrm{CVaR}_{.95}$ of held-out $\Delta$NLL} \\
    \midrule
GradDiff & \tblph & \tblph & \tblph & \tblph & \tblph & \tblph & \tblph & \tblph/5 & \tblph \\
RMU      & \tblph & \tblph & \tblph & \tblph & \tblph & \tblph & \tblph & \tblph/5 & \tblph \\
    \bottomrule
  \end{tabular}
\end{table*}

\subsection{The Sealed Score Improves Fixed-Budget Repair}
\label{sec:exp-protection}
\label{sec:exp-alpha}

\paragraph{Answer (RQ3).}
{\setlength{\emergencystretch}{3em}\ProtectionHeadline\par}

\paragraph{Selector comparison.}
Each confirmatory repair arm starts from the common criterion-reaching parent
checkpoint. They share the repair operator, neutral stream, example order,
optimizer random stream, token budget, guards, and save schedule; only the
$K_p$ selected discovery behaviors change. Panel C of
Figure~\ref{fig:profile-two-uses} summarizes this controlled comparison.

The confirmatory outcomes are mean and $\mathrm{CVaR}_{.95}$ damage on the
untouched audit fold. Lower values are better. On both core parents of the
primary setting, the two-signal arm attains the lowest held-out mean and
tail damage among the same-budget selectors while every constraint slack
stays positive (Table~\ref{tab:prot-effect}): the table reports our arm's
held-out damage beside the
strongest alternative selector and the unrepaired checkpoint, together with
the worst point estimates and worst individual one-sided 95\% bounds for paired
damage and native-metric non-inferiority, plus the minimum forgetting/utility
slack. This separates selector quality from
differences in repair compute or checkpoint admissibility.

\paragraph{Net repair benefit.}
Single-signal and random selectors test allocation quality under the same
repair capacity. The no-repair checkpoint answers a
different but necessary question: whether the repair pipeline provides net
benefit at all. A selector cannot pass RQ3 merely by winning among repaired
arms if repair itself worsens held-out behavior or violates the common
contract; the contract block of Table~\ref{tab:prot-effect} keeps this
condition beside the effect it licenses.

\paragraph{Diagnostics and boundaries.}
The $q_G^\star$-substituted mixture selector, uniform rehearsal, and the
realized-damage oracle remain secondary diagnostics outside RQ3.
Appendix~\ref{app:repair-protocol} specifies the shared operator, guards, and
rollback rule that make the arms comparable, and records which further
selection-quota and motion controls stay artifact-level rather than entering
the paper. Infeasible arms, incomplete random-draw
rosters, missing native outcomes, and forgetting or utility regressions remain
failures rather than being dropped.

% Coverage funnel over the executed campaign only. Rows are the four
% executed settings; every unfinished chain stays visible in Phase.
% NOTE: sections/generated/table_robustness.tex still emits the previous
% predeclared-roster layout and must be re-synced.

\begin{table}[!t]
  \caption{\textbf{Breadth and failure boundaries of the executed campaign.}
  Each row is one executed setting; \emph{Phase} is the furthest sealed
  phase reached (fidelity $\to$ calibration $\to$ freeze $\to$ audit $\to$
  mixture development $\to$ repair). RQ1 and RQ3 count claim-eligible and
  passing core parents after the pre-audit roster reduction; \emph{n.a.}
  marks a setting whose mixture weight and repair protocol are not frozen,
  so the corresponding claim is out of scope rather than failed. RQ2
  reports the registered setting-level fidelity certificate. Stress
  settings cannot rescue a primary failure.}
  \label{tab:robustness}
  \centering
  \small
  \setlength{\tabcolsep}{5pt}
  \begin{tabular}{@{}lcccc@{}}
    \toprule
    Setting & Phase reached & RQ1 parents & RQ2 cert. & RQ3 parents \\
    \midrule
TOFU / Qwen2.5-7B      & \tblph & \tblph/2 & pass & \tblph/2 \\
TOFU / Qwen2.5-1.5B    & \tblph & \textit{n.a.} & pass & \textit{n.a.} \\
TOFU / Llama-3.1-8B    & \tblph & \textit{n.a.} & pass & \textit{n.a.} \\
TOFU / Qwen2.5-14B     & \tblph & \textit{n.a.} & \tblph & \textit{n.a.} \\
WMDP-bio / Qwen2.5-7B  & \tblph & \textit{n.a.} & \tblph & \textit{n.a.} \\
WMDP-bio / Qwen2.5-14B & \tblph & \textit{n.a.} & \tblph & \textit{n.a.} \\
    \bottomrule
  \end{tabular}
\end{table}

\subsection{Breadth and Failure Boundaries}
\label{sec:exp-supporting}

\paragraph{Answer.}
{\setlength{\emergencystretch}{3em}\TransferHeadline\par}

Table~\ref{tab:robustness} keeps every executed model, benchmark, request,
and parent in view instead of pooling them into one universal effect. Each
row reports the furthest sealed phase a setting reached together with its
eligible and passing counts, so unfinished chains, invalid
profiles, parent non-reach, endpoint reversals, diffuse damage, infeasible
repair, and contract regressions remain visible. Breadth requires evidence in both output-
and representation-readout parent families and replication across the
executed model and benchmark settings. Readout labels describe objectives;
they never choose a signal endpoint or alter a target rule.


